% IEEE-style fragments (Abstract + Methodology) for: Execution-grounded
% hallucination mitigation in NL-to-SQL over Oracle.
% No fabricated numeric results — cite measured outcomes only from your own
% sql_evaluation_*.json runs when writing the final Results section.
%
% Usage: % IEEE-style fragments (Abstract + Methodology) for: Execution-grounded
% hallucination mitigation in NL-to-SQL over Oracle.
% No fabricated numeric results — cite measured outcomes only from your own
% sql_evaluation_*.json runs when writing the final Results section.
%
% Usage: % IEEE-style fragments (Abstract + Methodology) for: Execution-grounded
% hallucination mitigation in NL-to-SQL over Oracle.
% No fabricated numeric results — cite measured outcomes only from your own
% sql_evaluation_*.json runs when writing the final Results section.
%
% Usage: % IEEE-style fragments (Abstract + Methodology) for: Execution-grounded
% hallucination mitigation in NL-to-SQL over Oracle.
% No fabricated numeric results — cite measured outcomes only from your own
% sql_evaluation_*.json runs when writing the final Results section.
%
% Usage: \input{paper/ieee_abstract_methodology.tex}

\begin{abstract}
Natural-language interfaces to relational databases remain vulnerable to \emph{speculative} SQL: statements that appear plausible yet reference non-existent relations or columns, or that execute but answer the wrong question.
We describe an \emph{execution-grounded} pipeline for Oracle NL-to-SQL in which trust is anchored to observable database behavior rather than model prose alone.
A schema-aware prompt builder constrains generation to a documented table set; retrieval-augmented generation (RAG) supplies book excerpts and, when available, validated question--SQL pairs from an append-only query log.
Before presentation to the user, generated read-only SQL undergoes static checks against a machine-readable schema reference and optional on-server execution; the interface surfaces a structured verification verdict (static issues, parse/execute outcome, row counts, timing) and records explicit user feedback.
For offline benchmarking, an evaluation harness executes human-authored baseline SQL and model-generated SQL on the same Oracle instance, compares result sets under multiple equivalence notions, and records \texttt{EXPLAIN~PLAN} cost and latency deltas.
This design separates \emph{generation} from \emph{grounding}: hallucinations are detected when they contradict the schema, the executor, or (in batch evaluation) gold-intent result semantics---not when a language model merely sounds confident.
\end{abstract}

\section{Methodology}
\label{sec:methodology}

\subsection{System under study}
The studied system is a Node.js HTTP service that exposes NL-to-SQL generation (\texttt{POST /generate-sql}), optional Oracle execution via \texttt{oracledb}, and a browser lab UI.
Generation modes include lookup from curated question--SQL rules, LLM-only, and hybrid (lookup then LLM fallback).
The LLM is accessed through an OpenAI-compatible chat-completions endpoint; prompts embed fixed TPC-H-style table and column naming rules so that generated \texttt{SELECT} and \texttt{WITH\,\ldots\,SELECT} statements are intended to be copy-paste runnable on the lab schema.

\subsection{Retrieval-augmented generation (RAG)}
RAG combines three grounded sources: (i) a static \texttt{schema-reference.json} consumed by the UI and validation layer; (ii) optional full-text retrieval over an indexed course book EPUB (\texttt{book-index.js}) injected into generate and explain prompts when enabled; (iii) an optional query log (\texttt{data/query-log.jsonl}) of past \emph{verified} successful executions and user-supplied corrections, selected by lightweight lexical overlap with the current question and injected as short exemplar pairs.
RAG is treated as \emph{context}, not ground truth: the model is instructed to ignore off-topic passages and never to invent identifiers absent from the schema.

\subsection{Execution-grounded verification (online)}
For interactive use, the client may request \texttt{verify\_with\_oracle=true} on generation.
The server then (a) runs static analysis: referenced tables after \texttt{FROM}/\texttt{JOIN} are checked against the schema reference, with allowances for common Oracle builtins (e.g., \texttt{DUAL}) and for \texttt{WITH}-clause names; prefixed identifiers matching TPC-H column patterns are checked against the declared column set.
(b) If execution is enabled and credentials are present, the generated statement is executed once on Oracle (row cap aligned with the lab UI); parse and runtime errors, empty results, row counts, and elapsed time are returned in a structured \texttt{verification} object together with a qualitative risk summary synthesizing static and execution signals.
The UI may reuse returned rows to avoid duplicate round-trips when previews are included.
User thumbs-up/down and optional corrected SQL are appended to \texttt{data/sql-feedback.jsonl}; negative feedback with a correction is also written back to the query log to strengthen future retrieval.

\subsection{Batch evaluation against gold SQL (offline)}
Offline rigor uses a Python orchestrator (\texttt{experiments/run\_sql\_evaluation.py}) independent of the interactive path.
For each benchmark tuple (natural-language question, human-written SQL), the harness executes the baseline SQL on Oracle, then requests generated SQL from the HTTP API, executes it, and records success, latency, and \texttt{EXPLAIN~PLAN} metadata.
Equivalence of result sets is assessed with three complementary checks already implemented in the harness: multiset semantic equality (order ignored), exact row-and-order equality, and canonical string equality of normalized row tuples.
Disagreements and errors are exported for qualitative coding (\texttt{export\_failure\_cases.py}); aggregate metrics and plots are produced from the emitted JSON (\texttt{visualize\_results.py}).
\emph{All quantitative claims in the camera-ready paper must be computed from a named evaluation JSON artifact (timestamp, configuration, model ID); we do not state accuracy or cost figures here because they are run-specific.}

\subsection{Threats to validity and scope}
The static checker is syntactic and heuristic (it is not a full SQL semantic analyzer).
Online verification without gold SQL cannot prove \emph{intent} correctness---only schema consistency, executability, and coarse signals such as emptiness or runtime errors.
Batch semantic match requires gold SQL aligned with question intent; mismatches in the gold set limit conclusion strength.
Finally, LLM and Oracle versions, session parameters, and data skew must be fixed and reported for reproducibility.


\begin{abstract}
Natural-language interfaces to relational databases remain vulnerable to \emph{speculative} SQL: statements that appear plausible yet reference non-existent relations or columns, or that execute but answer the wrong question.
We describe an \emph{execution-grounded} pipeline for Oracle NL-to-SQL in which trust is anchored to observable database behavior rather than model prose alone.
A schema-aware prompt builder constrains generation to a documented table set; retrieval-augmented generation (RAG) supplies book excerpts and, when available, validated question--SQL pairs from an append-only query log.
Before presentation to the user, generated read-only SQL undergoes static checks against a machine-readable schema reference and optional on-server execution; the interface surfaces a structured verification verdict (static issues, parse/execute outcome, row counts, timing) and records explicit user feedback.
For offline benchmarking, an evaluation harness executes human-authored baseline SQL and model-generated SQL on the same Oracle instance, compares result sets under multiple equivalence notions, and records \texttt{EXPLAIN~PLAN} cost and latency deltas.
This design separates \emph{generation} from \emph{grounding}: hallucinations are detected when they contradict the schema, the executor, or (in batch evaluation) gold-intent result semantics---not when a language model merely sounds confident.
\end{abstract}

\section{Methodology}
\label{sec:methodology}

\subsection{System under study}
The studied system is a Node.js HTTP service that exposes NL-to-SQL generation (\texttt{POST /generate-sql}), optional Oracle execution via \texttt{oracledb}, and a browser lab UI.
Generation modes include lookup from curated question--SQL rules, LLM-only, and hybrid (lookup then LLM fallback).
The LLM is accessed through an OpenAI-compatible chat-completions endpoint; prompts embed fixed TPC-H-style table and column naming rules so that generated \texttt{SELECT} and \texttt{WITH\,\ldots\,SELECT} statements are intended to be copy-paste runnable on the lab schema.

\subsection{Retrieval-augmented generation (RAG)}
RAG combines three grounded sources: (i) a static \texttt{schema-reference.json} consumed by the UI and validation layer; (ii) optional full-text retrieval over an indexed course book EPUB (\texttt{book-index.js}) injected into generate and explain prompts when enabled; (iii) an optional query log (\texttt{data/query-log.jsonl}) of past \emph{verified} successful executions and user-supplied corrections, selected by lightweight lexical overlap with the current question and injected as short exemplar pairs.
RAG is treated as \emph{context}, not ground truth: the model is instructed to ignore off-topic passages and never to invent identifiers absent from the schema.

\subsection{Execution-grounded verification (online)}
For interactive use, the client may request \texttt{verify\_with\_oracle=true} on generation.
The server then (a) runs static analysis: referenced tables after \texttt{FROM}/\texttt{JOIN} are checked against the schema reference, with allowances for common Oracle builtins (e.g., \texttt{DUAL}) and for \texttt{WITH}-clause names; prefixed identifiers matching TPC-H column patterns are checked against the declared column set.
(b) If execution is enabled and credentials are present, the generated statement is executed once on Oracle (row cap aligned with the lab UI); parse and runtime errors, empty results, row counts, and elapsed time are returned in a structured \texttt{verification} object together with a qualitative risk summary synthesizing static and execution signals.
The UI may reuse returned rows to avoid duplicate round-trips when previews are included.
User thumbs-up/down and optional corrected SQL are appended to \texttt{data/sql-feedback.jsonl}; negative feedback with a correction is also written back to the query log to strengthen future retrieval.

\subsection{Batch evaluation against gold SQL (offline)}
Offline rigor uses a Python orchestrator (\texttt{experiments/run\_sql\_evaluation.py}) independent of the interactive path.
For each benchmark tuple (natural-language question, human-written SQL), the harness executes the baseline SQL on Oracle, then requests generated SQL from the HTTP API, executes it, and records success, latency, and \texttt{EXPLAIN~PLAN} metadata.
Equivalence of result sets is assessed with three complementary checks already implemented in the harness: multiset semantic equality (order ignored), exact row-and-order equality, and canonical string equality of normalized row tuples.
Disagreements and errors are exported for qualitative coding (\texttt{export\_failure\_cases.py}); aggregate metrics and plots are produced from the emitted JSON (\texttt{visualize\_results.py}).
\emph{All quantitative claims in the camera-ready paper must be computed from a named evaluation JSON artifact (timestamp, configuration, model ID); we do not state accuracy or cost figures here because they are run-specific.}

\subsection{Threats to validity and scope}
The static checker is syntactic and heuristic (it is not a full SQL semantic analyzer).
Online verification without gold SQL cannot prove \emph{intent} correctness---only schema consistency, executability, and coarse signals such as emptiness or runtime errors.
Batch semantic match requires gold SQL aligned with question intent; mismatches in the gold set limit conclusion strength.
Finally, LLM and Oracle versions, session parameters, and data skew must be fixed and reported for reproducibility.


\begin{abstract}
Natural-language interfaces to relational databases remain vulnerable to \emph{speculative} SQL: statements that appear plausible yet reference non-existent relations or columns, or that execute but answer the wrong question.
We describe an \emph{execution-grounded} pipeline for Oracle NL-to-SQL in which trust is anchored to observable database behavior rather than model prose alone.
A schema-aware prompt builder constrains generation to a documented table set; retrieval-augmented generation (RAG) supplies book excerpts and, when available, validated question--SQL pairs from an append-only query log.
Before presentation to the user, generated read-only SQL undergoes static checks against a machine-readable schema reference and optional on-server execution; the interface surfaces a structured verification verdict (static issues, parse/execute outcome, row counts, timing) and records explicit user feedback.
For offline benchmarking, an evaluation harness executes human-authored baseline SQL and model-generated SQL on the same Oracle instance, compares result sets under multiple equivalence notions, and records \texttt{EXPLAIN~PLAN} cost and latency deltas.
This design separates \emph{generation} from \emph{grounding}: hallucinations are detected when they contradict the schema, the executor, or (in batch evaluation) gold-intent result semantics---not when a language model merely sounds confident.
\end{abstract}

\section{Methodology}
\label{sec:methodology}

\subsection{System under study}
The studied system is a Node.js HTTP service that exposes NL-to-SQL generation (\texttt{POST /generate-sql}), optional Oracle execution via \texttt{oracledb}, and a browser lab UI.
Generation modes include lookup from curated question--SQL rules, LLM-only, and hybrid (lookup then LLM fallback).
The LLM is accessed through an OpenAI-compatible chat-completions endpoint; prompts embed fixed TPC-H-style table and column naming rules so that generated \texttt{SELECT} and \texttt{WITH\,\ldots\,SELECT} statements are intended to be copy-paste runnable on the lab schema.

\subsection{Retrieval-augmented generation (RAG)}
RAG combines three grounded sources: (i) a static \texttt{schema-reference.json} consumed by the UI and validation layer; (ii) optional full-text retrieval over an indexed course book EPUB (\texttt{book-index.js}) injected into generate and explain prompts when enabled; (iii) an optional query log (\texttt{data/query-log.jsonl}) of past \emph{verified} successful executions and user-supplied corrections, selected by lightweight lexical overlap with the current question and injected as short exemplar pairs.
RAG is treated as \emph{context}, not ground truth: the model is instructed to ignore off-topic passages and never to invent identifiers absent from the schema.

\subsection{Execution-grounded verification (online)}
For interactive use, the client may request \texttt{verify\_with\_oracle=true} on generation.
The server then (a) runs static analysis: referenced tables after \texttt{FROM}/\texttt{JOIN} are checked against the schema reference, with allowances for common Oracle builtins (e.g., \texttt{DUAL}) and for \texttt{WITH}-clause names; prefixed identifiers matching TPC-H column patterns are checked against the declared column set.
(b) If execution is enabled and credentials are present, the generated statement is executed once on Oracle (row cap aligned with the lab UI); parse and runtime errors, empty results, row counts, and elapsed time are returned in a structured \texttt{verification} object together with a qualitative risk summary synthesizing static and execution signals.
The UI may reuse returned rows to avoid duplicate round-trips when previews are included.
User thumbs-up/down and optional corrected SQL are appended to \texttt{data/sql-feedback.jsonl}; negative feedback with a correction is also written back to the query log to strengthen future retrieval.

\subsection{Batch evaluation against gold SQL (offline)}
Offline rigor uses a Python orchestrator (\texttt{experiments/run\_sql\_evaluation.py}) independent of the interactive path.
For each benchmark tuple (natural-language question, human-written SQL), the harness executes the baseline SQL on Oracle, then requests generated SQL from the HTTP API, executes it, and records success, latency, and \texttt{EXPLAIN~PLAN} metadata.
Equivalence of result sets is assessed with three complementary checks already implemented in the harness: multiset semantic equality (order ignored), exact row-and-order equality, and canonical string equality of normalized row tuples.
Disagreements and errors are exported for qualitative coding (\texttt{export\_failure\_cases.py}); aggregate metrics and plots are produced from the emitted JSON (\texttt{visualize\_results.py}).
\emph{All quantitative claims in the camera-ready paper must be computed from a named evaluation JSON artifact (timestamp, configuration, model ID); we do not state accuracy or cost figures here because they are run-specific.}

\subsection{Threats to validity and scope}
The static checker is syntactic and heuristic (it is not a full SQL semantic analyzer).
Online verification without gold SQL cannot prove \emph{intent} correctness---only schema consistency, executability, and coarse signals such as emptiness or runtime errors.
Batch semantic match requires gold SQL aligned with question intent; mismatches in the gold set limit conclusion strength.
Finally, LLM and Oracle versions, session parameters, and data skew must be fixed and reported for reproducibility.


\begin{abstract}
Natural-language interfaces to relational databases remain vulnerable to \emph{speculative} SQL: statements that appear plausible yet reference non-existent relations or columns, or that execute but answer the wrong question.
We describe an \emph{execution-grounded} pipeline for Oracle NL-to-SQL in which trust is anchored to observable database behavior rather than model prose alone.
A schema-aware prompt builder constrains generation to a documented table set; retrieval-augmented generation (RAG) supplies book excerpts and, when available, validated question--SQL pairs from an append-only query log.
Before presentation to the user, generated read-only SQL undergoes static checks against a machine-readable schema reference and optional on-server execution; the interface surfaces a structured verification verdict (static issues, parse/execute outcome, row counts, timing) and records explicit user feedback.
For offline benchmarking, an evaluation harness executes human-authored baseline SQL and model-generated SQL on the same Oracle instance, compares result sets under multiple equivalence notions, and records \texttt{EXPLAIN~PLAN} cost and latency deltas.
This design separates \emph{generation} from \emph{grounding}: hallucinations are detected when they contradict the schema, the executor, or (in batch evaluation) gold-intent result semantics---not when a language model merely sounds confident.
\end{abstract}

\section{Methodology}
\label{sec:methodology}

\subsection{System under study}
The studied system is a Node.js HTTP service that exposes NL-to-SQL generation (\texttt{POST /generate-sql}), optional Oracle execution via \texttt{oracledb}, and a browser lab UI.
Generation modes include lookup from curated question--SQL rules, LLM-only, and hybrid (lookup then LLM fallback).
The LLM is accessed through an OpenAI-compatible chat-completions endpoint; prompts embed fixed TPC-H-style table and column naming rules so that generated \texttt{SELECT} and \texttt{WITH\,\ldots\,SELECT} statements are intended to be copy-paste runnable on the lab schema.

\subsection{Retrieval-augmented generation (RAG)}
RAG combines three grounded sources: (i) a static \texttt{schema-reference.json} consumed by the UI and validation layer; (ii) optional full-text retrieval over an indexed course book EPUB (\texttt{book-index.js}) injected into generate and explain prompts when enabled; (iii) an optional query log (\texttt{data/query-log.jsonl}) of past \emph{verified} successful executions and user-supplied corrections, selected by lightweight lexical overlap with the current question and injected as short exemplar pairs.
RAG is treated as \emph{context}, not ground truth: the model is instructed to ignore off-topic passages and never to invent identifiers absent from the schema.

\subsection{Execution-grounded verification (online)}
For interactive use, the client may request \texttt{verify\_with\_oracle=true} on generation.
The server then (a) runs static analysis: referenced tables after \texttt{FROM}/\texttt{JOIN} are checked against the schema reference, with allowances for common Oracle builtins (e.g., \texttt{DUAL}) and for \texttt{WITH}-clause names; prefixed identifiers matching TPC-H column patterns are checked against the declared column set.
(b) If execution is enabled and credentials are present, the generated statement is executed once on Oracle (row cap aligned with the lab UI); parse and runtime errors, empty results, row counts, and elapsed time are returned in a structured \texttt{verification} object together with a qualitative risk summary synthesizing static and execution signals.
The UI may reuse returned rows to avoid duplicate round-trips when previews are included.
User thumbs-up/down and optional corrected SQL are appended to \texttt{data/sql-feedback.jsonl}; negative feedback with a correction is also written back to the query log to strengthen future retrieval.

\subsection{Batch evaluation against gold SQL (offline)}
Offline rigor uses a Python orchestrator (\texttt{experiments/run\_sql\_evaluation.py}) independent of the interactive path.
For each benchmark tuple (natural-language question, human-written SQL), the harness executes the baseline SQL on Oracle, then requests generated SQL from the HTTP API, executes it, and records success, latency, and \texttt{EXPLAIN~PLAN} metadata.
Equivalence of result sets is assessed with three complementary checks already implemented in the harness: multiset semantic equality (order ignored), exact row-and-order equality, and canonical string equality of normalized row tuples.
Disagreements and errors are exported for qualitative coding (\texttt{export\_failure\_cases.py}); aggregate metrics and plots are produced from the emitted JSON (\texttt{visualize\_results.py}).
\emph{All quantitative claims in the camera-ready paper must be computed from a named evaluation JSON artifact (timestamp, configuration, model ID); we do not state accuracy or cost figures here because they are run-specific.}

\subsection{Threats to validity and scope}
The static checker is syntactic and heuristic (it is not a full SQL semantic analyzer).
Online verification without gold SQL cannot prove \emph{intent} correctness---only schema consistency, executability, and coarse signals such as emptiness or runtime errors.
Batch semantic match requires gold SQL aligned with question intent; mismatches in the gold set limit conclusion strength.
Finally, LLM and Oracle versions, session parameters, and data skew must be fixed and reported for reproducibility.
